% One Britannica scan as extracted by Nemotron Parse 3.5. Filled by src/nemotron_challenge/stages/latex.py.
% The TITLE, BODY and SOURCE markers (each wrapped in two percent signs) are replaced. Compile with XeTeX (tectonic).
\documentclass[10pt]{article}
\usepackage[a4paper,margin=18mm]{geometry}
\usepackage{fontspec}
\usepackage{multicol}
\usepackage{graphicx}
\usepackage[hidelinks]{hyperref}
\setlength{\parindent}{0pt}
\setlength{\parskip}{4pt}
\setlength{\columnsep}{8mm}
\begin{document}
\section*{Historical source: OPTICS}
\begin{multicols}{2}
\subsection*{\textit{continued:} OPTICS}
the eye at a: fo that none of the rays that fall above \emph{h} on below i, can be reflected to the eye at o; and the dif. stance between \_6 and i is equal to half the length of the object \emph{AB.} Hence it appears, that if a man fees like whole image in a plane looking glats, the part of the glats that reflects his image melt be jull half as long and half as broad as himself, let him fland at any distance from it whatever; and that this image mull appear just as far behind the glafs as he is before it. Thus, the man AB (No. 30.) viewing himself in the plane mirrored CD, which is jull half as long as himself, fees his whole image as at \emph{EF}, behind the glafs exactly equal to his own fzz. For, a ray AC, proceeding from his eye at A, and falling perpendicularly upon the surface of the glafs at G, is re-flected back to his eye, in the frame line CaI : and the eye of this image will appear at B, in the frame line pro-duced to E, beyond the glafs. And a ray BD, wrong from his foot, and falling obliquely on the glafs at D, will be reflected as obliquely on the other side of the perpendicular AD), in the direction D,K, and the foot of this image will appear at Ft, in the direction of the re-fled ray AD, produced to F, where is is cut by the right line ACE, drawn parallel to the right line ACE. Jull the fans as if the glafs were taken away, and a real man flood at F, equal in five to the man (handing as B: for to his eyes at A, the eye of the other man at E would be few in the direction of the line ACE; and the foot of the man as F would be seen by the eye A, in the direrection of the line ADF. \emph{Of the} MAGIC LANTERN. ABCD (No. 31.) is a tin lantern, with a tube nklm fixed in the field of it. This tube conflicts of two joints, one of which dips into the others and by drawing this joint out, or pulling it in, the tube may be made longer or shorter. At kl in the end of the moveable joint of the tube, a convex lens is fixed, and an object painted with transparent colors upon a piece of thin glafs is plaedat de homewhere in the immovable joint of the tube; fo that as the tube is lengthened or fluorescent, the lens will be either at a greater or a left distance from this transferen object. In the field of the lantern three is a very convex lens bbc, which ferries to call a very frong light from the candle within the lantern upon the object de. New when the rays, which fringe through the objet de, diverge from the several points as Å, a, sc. in the object, and fall upon the lens kl, they will be made to converge to as many points fj, g \&e. on the other ride of the lens, and will point an inverted picture of the Vox. III. No §7. object at fg upon a white wall, a fiecet or a fireen of white paper, provided the object is farther from the less than its principal focus. To make this picture appear, dificent and brighty, it must have no other light fall upon it but what comes through the lens kl ; and for this reason the whole apparatus is to be placed in a dark, room EFGH. The lens A mull be very convex, fo that the object de may be very near to it, and yet not be nearorthan its principal focus; for by this means, as the object is near to the lens, the picture fg will be at a great distance from it , and consequently the picture will be much bigger than the object. Since the picture is inverted in respect of the object, in order to make the picture appear with the right end upwards, it is neceflary that the object de should be placed with the wrong end upwards. \emph{Of the} DIFFERENT REFRANGIBILITY of LIGHT. We have hitherto toppofed that a particle of light, as it comes from the fun, is the leaf particular iron which light can be separated. But we mull now correct this top-polition, by thawing, that a particle of light, as it comes from the fun, is, properly speaking, a bundle of rays, which may be farated from one another. Therefore, for the future, by a ray of light we mull be understood to mean, not that collection of particles which we have hitherto called by this name, but the leaf particles into which light can be farated. \_Rays of light are laid to be differently refrangible, when an the form or equal angles of incidence forms are more turned out of the way than others.\_ \_Rays are laid to be differently refasible, if firms are more safely reflected than others.\_ \_Light is called homogeneous, when all the rays are equality refrangible: i is treated heterogeneous, when time rays are other refrangible than others.\_ \_The colours of homogeneous light are called primary or female colours ; and thief of heterogeneous light are called secondary or mixed.\_ \_The rays of the fun are not all equally refrangible: and thief rays, which have a different degree of re-fragibility, here likewise a different colors.\_ If a beam of light SF, (No. 32.) that comes from the fun, paths into a dark room through F a round hole in a window fluter EG ; this beam proceeding straight forwards, and falling upon a paper as Y, would there make a round picture of the fun. This picture would be a confufed one indeed, if the hole is a large one and there is no lens in the hole. However, as this round f spotlight a a picture of the fun, we shall hereafter, notwithflanking its condition, call it by this name. Now it a glafs prifm ABC is placed between the hole in the window-fluiter and the paper at Y, the rays of this term; by the refraction which they differ in the prifm, will be bent from their fruit courte, and instead of going to io as to fall upon the paper at Y, they will be turned upwards, and the picture of the fun produced by them will fall upon a paper MN that is placed above Y. If all the rays
\end{multicols}
\section*{Source record}
\begin{itemize}
\item Page ID: \texttt{britannica:309f765b1ecc83e0a7f7762e44b4a927a7555443:train:681}
\item Scan: \texttt{page-0681-c2be06f2d615.jpg}, 600 x 855 px
\item Image SHA-256: \texttt{\small c2be06f2d6151658124b91589cab6163caa72e653bc6e1d62ecc30bbf9a202ee}
\item Extraction: \texttt{nvidia/nemotron-3.5-parse}, raw reply \texttt{artifacts/extract/raw/0681.json}
\item Text is Parse OCR output, normalized (long s as s, line-end hyphens joined). It is not corrected by hand.
\end{itemize}
\end{document}
